\section{Related Work}
\label{sec:related}
We situate our work at the intersection of three research areas.

\textbf{Agent evaluation and diagnosis.} Prior work on agent evaluation has focused on building comprehensive benchmarks~\citep{liu2023benchmarking,zheng2023judging} and measuring aggregate performance. We differ by diagnosing \emph{why} agents fail, building a structured taxonomy that enables per-category repair analysis. The closest prior work is \citet{zeng2024token}, which studies step-level credit assignment in DPO, and \citet{lai2024step} which proposes Step-DPO for redistributing preference weights to critical steps. Both focus on training-based credit assignment, while we focus on inference-time repair of structurally diagnosed failure categories.

\textbf{Training-free agent improvement.} Several works have proposed inference-time strategies for improving agents, including memory-augmented agents~\citep{le2022memory}, planner-executor splits~\citep{yao2023react}, and multi-agent collaboration~\citep{li2023multiagent}. We differ by systematically comparing these strategies against each other and against small-data SFT, using a shared taxonomy that clarifies which strategies repair which failure types. Prior comparisons~\citep{rporto2024agentrecipes} are conducted without a failure taxonomy and without controlled SFT baselines, making it difficult to identify the marginal value of each recipe.

\textbf{Failure analysis for LLMs.} Recent work has analyzed failure modes in LLMs for specific tasks such as math reasoning~\citep{tran2024understanding}, code generation~\citep{zhang2024debugging}, and tool use~\citep{patil2024evaluating}. We contribute a general failure taxonomy for multi-step agents that is benchmark-agnostic, with validation through a planned inter-annotator agreement study, and designed to be actionable for repair. Our taxonomy is related to but distinct from cognitive error taxonomies for humans~\citep{stanovich2011robot}.

\textbf{Small-data fine-tuning.} The question of how much data is needed for effective fine-tuning has been studied in domain adaptation~\citep{chen2021evaluate,zheng2023judging} and more recently in instruction tuning~\citep{zhou2024data}. We contribute a planned empirical study specifically focused on the interaction between SFT data scale and failure category, testing the hypothesis that category E (skill utilization deficiency) is the primary beneficiary of training and that the scaling curve has a threshold between 300 and 1,000 samples.
